\section{ResidueDroid Design}
\label{sec:design}

\subsection{Scope and threat model}
\label{sec:scope}

\sys is a \emph{cooperative} mechanism. The developer, or an SDK vendor, wants the app to keep its stated contract and links
\sys to do so. We do not defend against a malicious app, which can always exfiltrate data through a path we do not
mediate~\cite{reardon2019fiftyways}. The adversary is the schedule: a withdrawal arrives at an arbitrary time relative to queued
work, process death, reboot, force-stop and connectivity changes. A \emph{supported path} is a flow whose data is stored through
\sys's labelled stores and whose network submission goes through \sys's egress wrapper. Everything else is \emph{uncovered}.
That includes native code (JNI), custom network stacks such as raw sockets or embedded HTTP clients that bypass the wrapper,
WebViews, IPC to other apps, and any flow the developer did not route through the wrapper. We claim nothing for uncovered
paths. The raw-socket tenant measures one directly in \S\ref{sec:eval}.

\sys also cannot recall bytes. Data handed to the network before the withdrawal takes effect may still arrive afterwards, and
server-side copies are outside the device. The guarantee orders each accepted socket prefix relative to policy updates. A write can accept
only part of a request; its authorization log, local acceptance and remote receipt are distinct.

\subsection{Three mechanisms}
\begin{figure}[t]
\centering
\begin{tikzpicture}[font=\scriptsize,
  st/.style={draw=mSlate,rounded corners=2pt,fill=mNavy!7,align=center,inner sep=2.5pt,minimum height=0.7cm},
  op/.style={draw=mSlate,rounded corners=2pt,fill=mOchre!16,align=center,inner sep=2.5pt,minimum height=0.7cm},
  fe/.style={draw=mTeal,thick,rounded corners=2pt,fill=mTeal!11,align=center,inner sep=2.5pt,minimum height=0.85cm},
  arr/.style={-{Latex[length=1.5mm]},thick,draw=mInk!85}]
\node[op,text width=1.15cm] (col) at (0,0) {collect\\$\ell{=}\{(s,\epsilon)\}$};
\node[st,text width=1.3cm] (rows) at (1.8,0.5) {outbox rows\\+ label};
\node[st,text width=1.3cm] (der) at (1.8,-0.5) {derived rows\\$\ell{=}\bigcup$ inputs};
\node[st,text width=1.25cm] (job) at (3.55,0) {persisted job\\extras + label};
\node[fe,text width=1.2cm] (fence) at (5.55,0) {\icon[0.9em]{shield} fence\\check each\\chunk};
\node[draw=mSlate!70,rounded corners=2pt,fill=mSlate!8,align=center,text width=0.7cm,minimum height=0.7cm] (net) at (7.0,0) {\icon{cloud}\\socket};
\node[op,text width=1.55cm] (wd) at (3.55,-1.35) {withdraw (txn):\\$\epsilon{+}{+}$, $a{=}0$};
\node[st,text width=1.15cm] (gr) at (5.55,-1.35) {grants $\epsilon,a$\\(durable)};
\draw[arr] (col) -- (rows); \draw[arr] (col) -- (der); \draw[arr] (rows.east) -- (job); \draw[arr] (der.east) -- (job);
\draw[arr] (job) -- (fence);
\node[font=\tiny,align=center] at (4.3,0.72) {restored after kill / reboot};
\draw[arr] (fence) -- (net);
\draw[arr] (wd) -- (gr); \draw[arr,dashed] (gr) -- node[right,font=\tiny]{read} (fence);
\draw[thick,mTeal,dotted] (wd.north east) -- node[above left=-1pt,font=\tiny,text=mTeal]{lock $L$} (fence.south west);
\end{tikzpicture}
\caption{\sys data path. Labels travel with every persisted copy. The fence reads durable epochs and writes each bounded chunk while holding the lock that withdrawals also take.}
\Description{Source labels are stored with outbox rows, combined for derived rows and copied into persisted job extras. A fence checks each outgoing chunk against durable grants while sharing a lock with withdrawal.}\label{fig:design}
\end{figure}

Figure~\ref{fig:design} shows how the three mechanisms fit together.


\paragraph{Durable epochs.} For every withdrawable scope $s$, such as \textit{analytics} or \textit{profile}, \sys keeps a
durable epoch $\epsilon(s)$ and a flag $a(s)$. A withdrawal runs as one transaction that increments $\epsilon(s)$, clears
$a(s)$ and, for a deletion request, deletes local rows. A re-grant sets $a(s)$ again but keeps the epoch. Data collected
afterwards belongs to the new epoch, and data from older epochs stays invalid. The state lives in the app's SQLite database, so
it survives process death and reboot.

\paragraph{Persistent job lineage.} Each persisted datum carries a label $\ell(d)\subseteq\{(s,\epsilon)\}$, the scopes and
epochs of the sources it was computed from. The label is written \emph{in the same row or the same job input} as the data,
so it is committed atomically with the data. Collection assigns $\{(s,\epsilon(s))\}$, and derivation takes the union of its
inputs' labels. When data is copied into a job's persisted input (a \texttt{PersistableBundle} or WorkManager \texttt{Data}),
the label is copied with it. The label therefore survives in the job store, which a database deletion does not reach.

\paragraph{Late, serialized egress fence.} Every supported submission passes through
\texttt{Fence.submit(purpose, items)}. The audited HTTP adapter connects outside the lock,
encodes the request with a conservative union of its source labels, and writes chunks of at most
4\,KiB. Before each nonblocking write it acquires the same lock as withdrawal and rechecks
those labels. A zero-byte return is no handoff; a positive return records exactly the accepted
prefix. Waiting for socket readiness and reading the response occur outside the lock.
If policy changes after a partial send, the adapter closes the connection and suppresses the
remainder; previously accepted bytes remain outstanding. The contract rule for an item with label
$(s,e)$ and purpose $p$ is
\begin{equation}
\begin{aligned}
\mathit{forbid}(s,e,p)\;=\;{}&p\in R_s\;\wedge\;\big[(\pi_s\wedge e<\epsilon(s))\\
&\vee\;(e=\epsilon(s)\wedge\neg a(s))\big],
\end{aligned}
\label{eq:rule}
\end{equation}
where $R_s$ is the set of purposes the contract revokes and $\pi_s$ says whether it covers pending data
(Table~\ref{tab:contracts}). An item is forbidden if any element of its label is forbidden. For an OS-permission tenant
$R_s=\emptyset$, so prior data is never blocked. For a processing restriction, $R_s=\{\textit{personalisation}\}$, so backup
continues. For a ``stop collecting'' opt-out, $\pi_s$ is false, so queued data may still flush.

\begin{algorithm}[t]
\caption{Audited HTTP fence, with source labels on every retry}
\label{alg:fence}
\begin{algorithmic}[1]
\Require purpose $p$, encoded bytes $B$, durable source union $\ell$
\State Connect outside $L$; set operation deadline
\While{unsent bytes remain}
\State Wait for socket readiness outside $L$
\State Acquire $L$ within deadline
\State Read durable policy; reject unknown or forbidden labels
\State Attempt nonblocking write of at most 4\,KiB
\State Record returned accepted count; release $L$
\EndWhile
\State Read response outside $L$; retain source labels on retry
\end{algorithmic}
\end{algorithm}

\paragraph{Why this placement.} An early check can leave a window when an applicable withdrawal precedes an unprotected send.
Cancellation can miss re-enqueued or running work if it does not reach the transport. A check at enqueue misses a job queued before the withdrawal. A check at
job start misses everything the job prepared before the withdrawal and submits after it, including retries. Our declared-scope egress baseline misses a derived source when that source is invalidated
but the declared scope still passes. A boolean
consent flag without epochs lets old data through after a re-grant. \S\ref{sec:theory} turns each gap into a statement, and
\S\ref{sec:eval} measures each one.

\paragraph{Cost and composition.} One app-wide coordinator combines a local mutex with an
exclusive file lock. This simple design can briefly delay unrelated scopes, and storage or
scheduling delays can exceed nominal deadlines. Five-second operation deadlines are best-effort
failure bounds, not real-time guarantees. A failed or pending withdrawal is never acknowledged
as successful. Layered cancellation and early checks can reduce wasted work; the stated prefix-level guarantee uses a check for each final chunk. The old whole-request HTTP-stream wrapper is retained as a stream-wrapper
comparison, not used to justify the channel handoff boundary.
